% =========================================================
% How to Read an Analysis Exercise
% =========================================================
\paragraph{How to Read an Analysis Exercise Before Writing Anything.}

The preceding sections have given you the vocabulary: six exercise types,
twelve proof moves, four levels. This section shows how to put them
together in the five minutes before you write your first symbol.

\begin{tcolorbox}[colback=propbox, colframe=propborder, arc=2pt,
  left=6pt, right=6pt, top=4pt, bottom=4pt,
  title={\small\textbf{Pre-Writing Checklist}}, fonttitle=\small\bfseries]
\begin{enumerate}[itemsep=4pt]
  \item \textbf{Identify the type.}
    Which of the six types is this: Proof, Verification,
    Prove-or-Counterexample, Construction, Computation, Characterization?
    Write the type at the top of your scratch paper before anything else.

  \item \textbf{State the hypotheses and the goal.}
    In plain English, what are you assuming? What are you trying to show?
    Do not proceed until both are stated explicitly.

  \item \textbf{Unpack every definition in the goal.}
    If the goal involves a named concept --- convergence, compactness,
    closed, bounded --- write out its definition in full.
    The unpacked goal is your actual target.

  \item \textbf{Identify the key move.}
    Which of the twelve moves is most likely to bridge the hypothesis
    to the unpacked goal? Which level does that move belong to?
    Knowing the level tells you roughly how deep the proof will go and
    what other moves will likely appear alongside it.

  \item \textbf{Sketch the shape in words.}
    Write in plain language: ``I will choose $\varepsilon > 0$.
    I will use convergence of $\{x_n\}$ to find $N_1$. I will use
    convergence of $\{y_n\}$ to find $N_2$. I will take
    $N = \max(N_1, N_2)$. Then I will apply the triangle inequality.''
    Only after this sketch do you fill in the symbols.
\end{enumerate}
\end{tcolorbox}

\begin{remark}[The shape tells you what you need before you need it]
Step~5 is the most important and the most skipped. Students who jump
directly to symbols frequently paint themselves into corners: they
choose $N$ for one condition and then need it to satisfy a second
condition that their choice does not support.

If you write the shape of the proof in words first, you see
\emph{how many conditions $N$ must satisfy} before you commit to a
specific value. Then you choose $N$ to satisfy all of them simultaneously.
The tail argument tells you how: take the maximum of however many
separate $N_i$'s the conditions require.

The shape also tells you how many terms the triangle inequality will
produce and therefore what the right $\varepsilon$-splitting fractions
are. These are not things you discover mid-proof; they are visible from
the sketch.
\end{remark}

\begin{remark}[Common proofs and their shapes]
The following table summarizes the shape of the most common proof
patterns in a first analysis course, listed by the move sequence
and the levels they traverse.

\medskip
\begin{center}
\renewcommand{\arraystretch}{1.4}
\small
\begin{tabular}{p{0.28\textwidth} p{0.45\textwidth} p{0.15\textwidth}}
\toprule
\textbf{Goal} & \textbf{Move sequence} & \textbf{Levels} \\
\midrule
Convergent $\Rightarrow$ Cauchy &
  Unpack both definitions; triangle inequality; $\varepsilon$-splitting;
  tail argument ($\max$ of two $N$'s). & L1--L3 \\
Uniqueness of limits &
  Contradiction (assume two limits); triangle inequality; choose
  $\varepsilon = d(x,y)/2$; tail argument. & L2--L4 \\
Convergent $\Rightarrow$ bounded &
  Unpack convergence; choose $\varepsilon = 1$; tail argument
  (bound $n \geq N$); take $\max$ of first $N$ terms. & L1--L3 \\
Verify $d$ is a metric &
  List four conditions; check non-negativity, identity, symmetry directly;
  case split for triangle inequality. & L1--L2 \\
Compute diameter of $[a,b]$ &
  Write diameter as supremum; give upper bound; exhibit achieving pair. & L1--L2 \\
\bottomrule
\end{tabular}
\end{center}
\end{remark}

\begin{remark}[A note on templates]
For each of the six exercise types, there is a standard template
structure. These templates live in Appendix~A of the companion
\emph{Primer for Real Analysis Exercises}, which you will receive
at the start of Volume~III. The templates are starting-point
structures, not scripts to be memorized. Their purpose is to ensure
that every solution has all required parts and that no structural
element is accidentally omitted.

What matters at this stage is the vocabulary: six types, twelve moves,
four levels. You now have the map. Volume~III will fill in the terrain.
\end{remark}
